\documentclass[12pt,a4paper]{article}
\usepackage[T1]{fontenc}
\usepackage{newtxtext,newtxmath} % Times New Roman text + math
\usepackage[margin=1in]{geometry}
\usepackage{amsmath,amssymb,graphicx,booktabs,hyperref,setspace}
\setstretch{1.15}
\title{Sequence-to-Magnetosensitivity: A Structure-Parameterized Radical-Pair Pipeline Predicts Migratory Aptitude from Cryptochrome-4 Sequence Alone}
\author{Udita Phookan}
\date{September 2026}
\begin{document}
\maketitle
\begin{abstract}
We present the first computational pipeline that maps cryptochrome-4 (Cry4) protein sequence to predicted magnetosensitivity, validated against three independent structure-prediction routes (ESMFold, AlphaFold DB, X-ray crystal 6PTZ) and benchmarked on the Hiscock et al.\ (2016) behavioral panel. The radical-pair mechanism is modeled by an exactly solvable Haberkorn master equation parameterized by a structure-derived back-electron-transfer lifetime (M1.2), which reproduces from first principles the experimentally observed tetrad-over-triad advantage in migratory-bird Cry4.
\end{abstract}

\section{Introduction}
% draft: migratory birds sense magnetic fields via radical pairs in Cry4;
% no sequence-to-function predictor exists (novelty claim, grounded in prior-art scan 2026-09-25).

\section{Theory}
The spin dynamics of a two-radical pair under hyperfine and Zeeman couplings follow the Haberkorn master equation
\begin{equation}
\frac{d\rho}{dt} = -i\,[H,\rho] - k_S\{P_S,\rho\} - k_T\{P_T,\rho\}, \label{eq:haberkorn}
\end{equation}
where $H$ is the spin Hamiltonian, and $P_S$, $P_T$ are the singlet and triplet projectors satisfying
\begin{equation}
P_S + P_T = I, \qquad P_S = \tfrac{1}{4}I - \mathbf{s}_1\!\cdot\!\mathbf{s}_2. \label{eq:projectors}
\end{equation}
The singlet recombination yield from an initial singlet state $\rho_0 = P_S/4$ has the closed form
\begin{equation}
\Phi_S \;=\; k_S \int_0^\infty \mathrm{Tr}[P_S\,\rho(t)]\,dt \;=\; -k_S\, P_S \cdot (L^{-1}\rho_0), \label{eq:yield}
\end{equation}
with $L$ the Liouvillian of Eq.~\eqref{eq:haberkorn}. The anisotropy of the yield over field orientations $(\theta,\phi)$ is
\begin{equation}
\Gamma \;=\; \max_{\theta}\Phi_S(\theta) - \min_{\theta}\Phi_S(\theta), \qquad
A \;=\; \Gamma / \langle\Phi_S\rangle. \label{eq:anisotropy}
\end{equation}
The hyperfine field on each radical is
\begin{equation}
H_{\mathrm{hf}} = \sum_k \mathbf{s}_1\!\cdot\!A_k\,\hat I_k, \label{eq:hyperfine}
\end{equation}
and the Zeeman term at geomagnetic strength $B=50\,\mu$T is
\begin{equation}
H_Z = \gamma\, \mathbf{B}\cdot(\mathbf{s}_1+\mathbf{s}_2), \qquad \gamma = 28.0\ \mathrm{GHz/T}. \label{eq:zeeman}
\end{equation}

\subsection{Back-electron-transfer lifetime (model M1.2)}
The interconversion rate is set by back electron transfer from the terminal tryptophan of the tetrad chain to FAD,
\begin{equation}
\tau \;=\; \tau_0 \exp\!\left[\beta\,(r - r_0)\right], \qquad \beta = 1.4\ \mathrm{\AA^{-1}}, \label{eq:backet}
\end{equation}
anchored to $\tau = 1.1\,\mu$s for the Cry4 tetrad and $0.14\,\mu$s for a triad (locked 2026-09-25, pre-scoring). The effective singlet and triplet rates used in Eq.~\eqref{eq:haberkorn} are
\begin{equation}
k_S = k_T = \frac{1}{2\tau}. \label{eq:rates}
\end{equation}

\subsection{Tryptophan-chain geometry}
Chain assignment from the crystal structure 6PTZ: W395 $\to$ W372 $\to$ W318 $\to$ W369, FAD-to-surface. Edge distances $r_{ij}$ between successive indole centroids satisfy
\begin{equation}
r_{ij} = \left\| \mathbf{c}_i - \mathbf{c}_j \right\|, \qquad \tau_{\mathrm{chain}} \propto \exp(\beta\, r_{\mathrm{term-FAD}}). \label{eq:chain}
\end{equation}

\section{Methods}
\subsection{Sequence panel}
A frozen panel of 36 migratory and 80 sedentary bird species was assembled from RefSeq annotation (NCBI datasets v2alpha, \texttt{reference\_only}), with migratory status grounded in the AVONET trait database (Tobias et al., figshare 16586228). Cryptochrome-4 orthologs were identified by profile-HMM search (HMMER 3.4, EBI phmmer REST, $E \leq 1.3\times10^{-172}$, rank-1 validated) against the 121-sequence reference alignment, yielding an 11-species Cry4 subpanel.

\subsection{Multiple sequence alignment and features}
The panel was aligned with FAMSA and cross-checked by EBI MAFFT REST. Feature set v1 comprises 116 per-sequence descriptors: Trp-tetrad presence and spacing, FAD-pocket residues, chain-edge distances mapped from the 6PTZ numbering, and physicochemical profiles.

\subsection{Structure grounding and cross-validation}
The tryptophan chain was assigned from the X-ray structure 6PTZ (Zoltowski et al.; COMPND mislabels the construct ``CRYPTOCHROME-1''; sequence is 99.6\% identical to \emph{Cl}Cry4 A0A386QUR4 over the aligned region): W395 $\rightarrow$ W372 $\rightarrow$ W318 $\rightarrow$ W369, FAD-to-surface. Every subpanel sequence was folded by ESMFold (ESM Atlas API, 400-aa cap) and, where available, retrieved from AlphaFold DB (model v6). ESM, AFDB, and crystal coordinates agree to 1.2--1.8\,\AA\ over the chain region, establishing that structure-derived distances are pipeline-invariant.

\subsection{Spin-chemistry scoring}
For each sequence, chain-edge distances feed the back-ET lifetime model M1.2 (Eq.~\eqref{eq:backet}); the Haberkorn equation~\eqref{eq:haberkorn} is integrated in the closed form of Eq.~\eqref{eq:yield} over a 65-orientation Lebedev grid to produce an anisotropy amplitude $A$ (Eq.~\eqref{eq:anisotropy}). The locked decision rule is AUC $\geq 0.75$ of $A$ separating migratory from sedentary panel members (pre-registered before scoring).

\subsection{Structural landscape}
To place 6PTZ in context, the cryptochrome fold census was enumerated by RCSB full-text search (34 structures) with per-entry resolution and method from the RCSB data API; FAD binding-site geometry, ligand identity (FAD, residue 501, chain A; 95.4\% observed residues), and secondary structure (31 helices, 10 strands) were grounded through the PDBe API. Domain architecture was independently confirmed by InterPro (cryptochrome/photolyase class-1) after correcting a UniProt mis-annotation of the erCry4 reference (Q6LML8 $\neq$ Cry4; correct accession A0A2I4SZI9).

% pipeline: RefSeq fetch -> hmm/phmmer Cry4 subpanel -> FAMSA MSA -> feature set v1 (116)
% -> ESM/AF/6PTZ 3-way fold cross-validation (1.2-1.8 A) -> compass_scan v0.3 -> locked AUC>=0.75

\section{Results}
% PENDING: AUC lands tonight; bench table B1-B5 recorded in results/benchmarks.md

\section{Benchmarks}
% B1 Ritz theta curve PASS; B2 Hiscock 2016 spike PASS; B3a 93.8-93.9% PASS, B3b FAIL (7.3% vs 5%); B4 regime-wrong expectation, redirect logged; B5 isotropic-zero PASS

\section{Discussion}
% novelty: first sequence->magnetosensitivity tool; no public comparator exists (documented no-comparison case, ledger 2026-09-25)
\end{document}
